\section{Conclusion}
\label{sec:conclusion}

In conclusion, the evidence indicates that the construction of federal prisons significantly altered the dynamics of Mexico’s criminal justice system. While prison expansion alleviated overcrowding in the long run, it also generated short-term increases in processing and sentencing, suggesting that capacity shocks reshape both caseload pressures and judicial outcomes. Importantly, the divergent trends across sentencing types—higher rates of incarceration but no change in fines—point to behavioral adjustments by judges rather than mechanical spillovers from greater caseloads alone. These findings underscore that judicial decision-making responds not only to legal frameworks but also to the institutional environment, particularly the availability of resources and prison capacity.

Overall, the results show that new prison construction has lasting effects beyond simply expanding space. It alters how judges sentence, how often pretrial detention is used, and how punishments are distributed. While sentence length remains largely unaffected, the findings demonstrate that judicial choices shift in response to capacity changes. By combining detailed judicial data with novel measures of prison infrastructure, this study provides rare causal evidence that investments in prisons can directly shape judicial behavior and outcomes in the criminal justice system.